% ============================================================================
\section{What it TAUght us}
\label{sec:lessons}
% ============================================================================

\paragraph{Design the loss for the decision.}
Building a task-specific loss meant translating the trigger into rewards, protected regions and a fake-rate constraint, rather than merely choosing a standard classification loss. Its value was control over the efficiency profile, not a guaranteed global gain (\cref{sec:loss,app:analysis_loss}).

\paragraph{More gradient is not more learnability.}
We initially up-weighted difficult low-energy taus to improve their efficiency; this often hurt. Emphasising cleaner, higher-energy examples worked better for several inputs. Our working explanation is that useful class structure can transfer, whereas ambiguous examples may consume gradient without improving separation. Weighting must reflect what the data and model can learn, not only where we want improvement; this is a problem-dependent lesson (\cref{app:weighting,app:analysis_weighting}).

\paragraph{Understand the input before enlarging the model.}
The earlier MLP sweep tested networks up to roughly $100$ times larger without improving its small reference, while our gains came from fine-grained EM2 information and compactness features. We learned to question input discrimination before blaming capacity: a larger network cannot restore detail absent from its inputs (\cref{sec:features,app:analysis_capacity,app:analysis_representation}).

\paragraph{Complexity must earn its parameters.}
Deeper or strided EM2 branches did not systematically beat the simple E1. Local patterns and their positions may already suffice on these small images; E1's dense head also combines them globally, while further compression can lose detail. This is a plausible explanation, not proof of why a particular model lost. We therefore favour the smallest architecture that preserves the required efficiency profile (\cref{sec:arch,app:analysis_capacity}).

\paragraph{Loss (function), but not for nothing.}
We learned three design requirements the hard way: (i)~reward better ranking, not score inflation; (ii)~give every relevant energy region an explicit reward or protection, or it can be sacrificed; and (iii)~ensure gradients reach missed signal, not just objects already near the threshold. Exact out-of-sample checks and statistically meaningful tolerances complete the recipe. \Cref{app:loss_requirements} details all five lessons and the failures that revealed them.

\paragraph{Pair learning follows the event decision.}
Average losses within eligible events and keep the selected systems fixed for
evaluation (\cref{app:pair_protocol}).
